% P15.32: the full Durrett Section 4.7 application and exercise chain.
\chapter{Exchangeability and de Finetti Representation}
\label{sec:p15-exchangeability-and-definetti}

Symmetry of an infinite sequence leads to conditional independence
given a random probability law. The conditioning is essential: the
unconditioned observations can remain dependent. This chapter supplies
the reverse-martingale applications and all exercises in Durrett's
Section~4.7, including its separate Hewitt--Savage proof, ballot
argument and de Finetti construction
\cite[pp.~249--255]{DurrettPTEAuthor2019}.

\section{Finite permutations and conditional averages}

\begin{hbdefinition}{def:p15-exchangeable-canonical}{The canonical law and decreasing invariant information}
Let \((S,\mathcal S)\) be a nonempty measurable space,
\(\Omega=S^{\mathbb N_+}\), \(\mathcal F=\mathcal S^{\otimes\mathbb N_+}\),
and \(X_j(\omega)=\omega_j\). Let \(P\) be the sequence law.
Retain the finite permutations \(T_\pi\) of
\cref{def:p15-permutable-events}. For integers \(n\geq1\), define
\[
 \mathcal E_n=\{A\in\mathcal F:
       T_\pi^{-1}A=A\text{ for every permutation of }1,\ldots,n\},
 \qquad \mathcal E=\bigcap_{n\geq1}\mathcal E_n.
\]
Each permutation in this formula fixes every coordinate beyond \(n\).
The tail sigma-field is
\(\mathcal T=\bigcap_{m\geq1}\sigma(X_j:j>m)\).
These are exact invariant-event definitions.

The law is exchangeable if each finite coordinate vector has the same
distribution after every permutation of its coordinates.
For \(n\geq k\geq1\), let \(I_{n,k}\) be the ordered tuples of
\(k\) distinct indices in \(\{1,\ldots,n\}\).
For a real measurable function \(\varphi:S^k\to\mathbb R\), define
\[
 (n)_k=n(n-1)\cdots(n-k+1),\qquad
 A_{n,k}(\varphi)=\frac1{(n)_k}
       \sum_{i\in I_{n,k}}\varphi(X_{i_1},\ldots,X_{i_k}).
\]
The falling-product notation is local to this chapter.
\index{exchangeable sequence}\index{symmetrization}
\end{hbdefinition}

\begin{hblemma}{thm:p15-exchangeable-canonical}{proved=annex:p15.exchangeable-canonical}{Measure invariance and the canonical reduction}
Finite-vector exchangeability is equivalent to preservation of \(P\)
by every \(T_\pi\) on all of \(\mathcal F\).
Each \(\mathcal E_n\) is a sigma-field,
\(\mathcal E_{n+1}\subseteq\mathcal E_n\), and
\(\mathcal T\subseteq\mathcal E\).
For a sequence on another probability space, its coordinate map
pushes that law to \(P\); the joint-law and all-event conditional
identities below pull back through that map. They concern the sequence
and its invariant information.
\end{hblemma}

\begin{hbtheorem}{thm:p15-exchangeable-averages}{proved=annex:p15.exchangeable-averages}{Every integrable finite symmetrization is a reverse conditional expectation}
Under an exchangeable law, let
\(\varphi:S^k\to\mathbb R\) be measurable with
\(E|\varphi(X_1,\ldots,X_k)|<\infty\).
For every \(n\geq k\),
\[
 A_{n,k}(\varphi)=E(\varphi(X_1,\ldots,X_k)\mid\mathcal E_n).
\]
Consequently these averages converge almost surely and in \(L^1\) to
\(E(\varphi(X_1,\ldots,X_k)\mid\mathcal E)\).
Different integrands may have different exceptional null sets.
\end{hbtheorem}

\begin{hblemma}{thm:p15-independent-conditional-projection}{proved=annex:p15.independent-conditional-projection}{An independent version of a conditional mean is constant}
Let \(V\in L^2\), \(Y=E(V\mid\mathcal G)\), and suppose \(Y\)
has a version measurable with respect to a sigma-field \(\mathcal H\)
independent of \(V\). Then \(Y=EV\) almost surely.
The sigma-fields \(\mathcal G,\mathcal H\) need not be nested.
\end{hblemma}

\begin{hbtheorem}{thm:p15-reverse-hewitt-savage}{proved=annex:p15.reverse-hewitt-savage}{The source's reverse-average proof of Hewitt--Savage}
For i.i.d. coordinates and bounded measurable \(\varphi:S^k\to\mathbb R\),
\[
                A_{n,k}(\varphi)\longrightarrow
                     E\varphi(X_1,\ldots,X_k)\quad\text{almost surely}.
\]
The proof deletes finitely many coordinates from the average and
uses the independent-conditional-projection lemma.
It then proves that every \(A\in\mathcal E\) has probability zero or one.
This is the distinct source route of Example~4.7.6 and Lemma~4.7.7;
it assumes neither a strong law nor the earlier independent-block
Hewitt--Savage proof.
\end{hbtheorem}

\section{The reverse strong law and the ballot calculation}

\begin{hbtheorem}{thm:p15-reverse-strong-law}{proved=annex:p15.reverse-strong-law}{The source's third strong-law route}
For i.i.d. real \(X_j\) with \(E|X_1|<\infty\), let
\(S_n=\sum_{j=1}^nX_j\), and
\[
 \mathcal H_n=\sigma(S_n,S_{n+1},\ldots)
             =\sigma(S_n,X_{n+1},X_{n+2},\ldots).
\]
The sigma-fields decrease, \(\mathcal H_n\subseteq\mathcal E_n\),
and \(S_n/n=E(X_1\mid\mathcal H_n)\).
The reverse conditional-limit theorem therefore gives
\(S_n/n\to EX_1\) almost surely and in \(L^1\).
This preserves the source Example~4.7.4 argument.
Its limiting cases (an infinite first moment, and an exchangeable but
not independent law whose limit is random) are worked in
\cref{thm:p15-w02b-p15-022-reverse-strong-law-boundary}.
\end{hbtheorem}

\begin{hbtheorem}{thm:p15-reverse-ballot}{proved=annex:p15.reverse-ballot}{The full nonnegative-integer ballot identity}
Let \(\xi_1,\ldots,\xi_N\) be a finite exchangeable vector of
nonnegative integer-valued variables, \(N\geq1\).
Put \(S_j=\sum_{i=1}^j\xi_i\) and
\(G=\{S_j<j\text{ for every }1\leq j\leq N\}\).
No moment assumption is imposed. Then
\[
                        P(G\mid S_N)=(1-S_N/N)^+ .
\]
This includes the source's i.i.d. hypothesis. Its reverse stopping
proof is applied under each positive-probability conditioning event
\(S_N=m<N\), where every coordinate is bounded.

In the election interpretation, independent fair values \(0,2\)
represent votes for candidates A and B. Given integer totals
\(\alpha>\beta\geq0\), \(N=\alpha+\beta>0\), the probability that
A leads strictly throughout counting is
\((\alpha-\beta)/(\alpha+\beta)\).

Nonnegativity and integer values matter. For \(N=1,\xi_1=1/2\),
the event has probability one but the proposed right side is \(1/2\).
For two independent fair values \(-1,1\), conditional on \(S_2=0\),
the event has probability \(1/2\) but the proposed right side is one.
\end{hbtheorem}

\section{Conditional independence and a random law}

\begin{hbtheorem}{thm:p15-definetti-factorization}{proved=annex:p15.definetti-factorization}{de Finetti's conditional-product identity}
Under an exchangeable law, for any \(k\geq1\) and bounded measurable
\(f_1,\ldots,f_k:S\to\mathbb R\),
\[
 E\!\left(\prod_{j=1}^kf_j(X_j)\,\middle|\,\mathcal E\right)
                   =\prod_{j=1}^k E(f_j(X_1)\mid\mathcal E)
                                      \quad\text{almost surely}.
\]
This is conditional independent identical distribution in the
all-event conditional-expectation sense.
The proof retains the source's exact distinct-index and repeated-index
decomposition. It applies on an arbitrary measurable coordinate space;
the probability-valued representation below has an additional
standard Borel hypothesis
\cite[thm.~4.7.9]{DurrettPTEAuthor2019}.
\end{hbtheorem}

\begin{hbdefinition}{def:p15-directing-law}{The measurable space of probability laws}
For a nonempty standard Borel space \((S,\mathcal S)\), let
\(\mathscr P(S)\) be its probability measures, with sigma-field
generated by all evaluations \(e_B(\mu)=\mu(B)\), \(B\in\mathcal S\).
Choose a countable determining pi-system \(\mathcal C\) containing \(S\)
and generating \(\mathcal S\); one is obtained from a Borel real coding
and rational half-lines.
A probability kernel \(K\) from \((\Omega,\mathcal E)\) to \(S\)
is thus a measurable probability-valued map \(\omega\mapsto K_\omega\).
Its distribution on \(\mathscr P(S)\) is denoted \(\Lambda\).
Write \(\sigma(K)=\sigma(K(B):B\in\mathcal S)\) for the information
generated by this random law.

For a sub-sigma-field \(\mathcal G\subseteq\mathcal F\), define
\[
 \mathcal G^{P}=\{A\in\mathcal F:
                 P(A\mathbin{\triangle}B)=0\text{ for some }B\in\mathcal G\}.
\]
This records equality of information modulo ambient measurable null events.
\index{directing random measure}\index{mixture!independent sequences}
\end{hbdefinition}

\begin{hbtheorem}{thm:p15-definetti-kernel}{proved=annex:p15.definetti-kernel}{The standard Borel mixture, on all sequence events}
Let \(S\) be standard Borel and \(P\) exchangeable.
Choose a regular conditional law \(K_\omega\) of \(X_1\) given
\(\mathcal E\). The map
\[
                         J(\omega,\cdot)=K_\omega^{\otimes\mathbb N_+}
\]
is a probability kernel and a regular conditional law of the complete
coordinate sequence given \(\mathcal E\).
Its mixing law \(\Lambda\) satisfies
\[
                  P(A)=\int_{\mathscr P(S)}
                          \mu^{\otimes\mathbb N_+}(A)\,\Lambda(d\mu)
                                      \quad(A\in\mathcal F).
\]
In particular for \(B_1,\ldots,B_k\in\mathcal S\),
\[
 P(X_1\in B_1,\ldots,X_k\in B_k)
           =\int\prod_{j=1}^k\mu(B_j)\,\Lambda(d\mu).
\]
The required kernel measurability and common countable determining
class are proved in the annex.
\end{hbtheorem}

\begin{hbtheorem}{thm:p15-directing-law-uniqueness}{proved=annex:p15.directing-law-uniqueness}{Recovery, uniqueness and invariant information}
For every \(C\in\mathcal C\), simultaneously almost surely,
\[
                   \frac1n\sum_{j=1}^n\mathbf1_{\{X_j\in C\}}
                                 \longrightarrow K_\omega(C).
\]
The law \(\Lambda\) in the preceding theorem is unique.
Also
\[
                      \mathcal E^{P}=\mathcal T^{P}
                                   =\sigma(K)^{P}.
\]
This is equality modulo null events, not literal equality of the
unaugmented sigma-fields.
The exchangeable sequence is i.i.d. exactly when its directing law
\(\Lambda\) is a point mass, or equivalently when
\(\mathcal E\) is probability-trivial.
\end{hbtheorem}

\begin{hbtheorem}{thm:p15-exchangeable-strong-laws}{proved=annex:p15.exchangeable-strong-laws}{Averages and empirical laws under exchangeability}
For every fixed real measurable \(f\) with \(E|f(X_1)|<\infty\),
\[
       \frac1n\sum_{j=1}^nf(X_j)\longrightarrow
                   E(f(X_1)\mid\mathcal E)
                          \quad\text{almost surely and in }L^1.
\]
In the standard Borel case the limit is \(\int f\,dK_\omega\).
If only \(\int|f|\,d\mu<\infty\) for \(\Lambda\)-almost every
\(\mu\) is assumed, the same almost-sure conclusion holds, without
an unconditional \(L^1\) assertion.

For real coordinates, the empirical probabilities
\(\widehat\mu_n=n^{-1}\sum_{j=1}^n\delta_{X_j}\)
converge weakly to \(K_\omega\) almost surely.
The proof uses one countable rational determining family before
passing to each limiting CDF continuity point.
\end{hbtheorem}

\section{Binary sequences and finite sampling}

\begin{hbtheorem}{thm:p15-binary-definetti}{proved=annex:p15.binary-definetti}{The binary mixture and the exact conditional sampling law}
For an infinite exchangeable binary sequence, there is a unique
probability measure \(\lambda\) on \([0,1]\) such that for every
word \(b\) of length \(n\) and count \(s\),
\[
             P(X_1=b_1,\ldots,X_n=b_n)
                    =\int_0^1\theta^s(1-\theta)^{n-s}\,\lambda(d\theta).
\]
The powers with exponent zero are one, including at endpoints.
The measure \(\lambda\) is the law of the almost-sure limiting
success frequency. Count probabilities have the additional factor
\(\binom ns\).

For a finite exchangeable binary vector of length \(N\), condition
only at \(m\) with \(P(S_N=m)>0\), where \(S_N=\sum_{j=1}^NX_j\).
All length-\(N\) words of that count are then equally likely, and
\[
              P(X_1=\cdots=X_k=1\mid S_N=m)
                  =\frac{\binom{N-k}{m-k}}{\binom Nm},
                           \quad 1\leq k\leq N.
\]
A binomial coefficient is zero when its lower integer lies outside
zero through its upper integer. This supplies Theorem~4.7.10 and
Exercise~4.7.3 at their actual infinite and finite scopes.
\end{hbtheorem}

\begin{hbtheorem}{thm:p15-finite-urn-approximation}{proved=annex:p15.finite-urn-approximation}{The finite-urn approximation and a second binary existence proof}
Let \(H_{m,N}^{(k)}\) be the law of the first \(k\) ordered draws
without replacement from \(N\) binary entries containing \(m\) ones,
where \(0\leq m\leq N\) and \(1\leq k\leq N\).
For the event-supremum total variation distance,
\[
 \left\|H_{m,N}^{(k)}
      -\operatorname{Bernoulli}(m/N)^{\otimes k}\right\|_{\rm TV}
                         \leq\frac{k(k-1)}{2N}.
\]
For any finite exchangeable binary vector, mix over its total
success count to get the same bound against the mixture whose
parameter law is that of \(S_N/N\).
For an infinitely exchangeable sequence, compact subsequence
selection of those parameter laws gives the binary representation
without using the general conditional-product theorem.
This completes the finite-urn argument described in the source's
post-Theorem~4.7.10 discussion. The displayed coupling bound is
proved here; it is not claimed to be a sharp bound from the
separately cited Diaconis--Freedman paper.
\end{hbtheorem}

\section{Moment constraints and the reverse-limit exercises}

\begin{hbtheorem}{thm:p15-exchangeable-second-moments}{proved=annex:p15.exchangeable-second-moments}{Positive common covariance and the pairwise variance statistic}
For an infinite exchangeable real sequence with \(EX_1^2<\infty\), put
\(M=E(X_1\mid\mathcal E)\). Then
\[
 E(X_1X_2)=EM^2\geq0,\qquad
 \operatorname{Cov}(X_1,X_2)=\operatorname{Var}(M)\geq0.
\]
For i.i.d. real \(X_j\) of variance \(\sigma^2<\infty\),
\[
            \binom n2^{-1}\sum_{1\leq i<j\leq n}(X_i-X_j)^2
                      \longrightarrow2\sigma^2
                            \quad\text{almost surely and in }L^1.
\]
The proof uses the integrable symmetrization identity, including
its unbounded integrand, as required by Exercises~4.7.4--4.7.5.

The finite exchangeable pair assigning probability \(1/2\) to each
of \((1,0),(0,1)\) has covariance \(-1/4\).
It admits neither an infinite exchangeable extension nor a mixture
of identical independent Bernoulli laws.
\end{hbtheorem}

\begin{hbtheorem}{thm:p15-backward-lp-limit}{proved=annex:p15.backward-lp-limit}{The finite-Lp reverse convergence exercise}
For decreasing \(\mathcal G_n\), an integrable \(Z\), and
\(Y_n=E(Z\mid\mathcal G_n)\), the limit is
\(Y_\infty=E(Z\mid\bigcap_n\mathcal G_n)\).
If \(Z\in L^p\) for finite \(1<p<\infty\), convergence also holds
in \(L^p\). For \(0<p\leq1\), the already known \(L^1\) convergence
implies \(L^p\) convergence. This supplies Exercise~4.7.1.

Boundedness alone does not give \(L^\infty\) convergence.
On \((0,1]\) with Lebesgue probability, put
\(a_n=1/(n+2)\), \(Z=\mathbf1_{(0,1/2]}\), and
\[
 \mathcal G_n=\sigma((0,1/(m+2)]:m\geq n).
\]
Then
\[
 Y_n=\mathbf1_{(0,a_n]}
       +\frac{1/2-a_n}{1-a_n}\mathbf1_{(a_n,1]},
 \qquad Y_\infty=1/2,
 \qquad \|Y_n-Y_\infty\|_\infty=1/2
\]
at every \(n\geq0\).
\end{hbtheorem}

\begin{hbtheorem}{thm:p15-backward-changing-variables}{proved=annex:p15.backward-changing-variables}{Changing the integrand along decreasing information}
Let \(\mathcal G_n\) decrease to \(\mathcal G_\infty=\bigcap_n\mathcal G_n\).
If \(Z_n\to Z\) almost surely and \(|Z_n|\leq W\) for every \(n\)
with \(W\geq0\) integrable, then
\[
                   E(Z_n\mid\mathcal G_n)
                      \longrightarrow E(Z\mid\mathcal G_\infty)
                         \quad\text{almost surely and in }L^1.
\]
The proof uses the decreasing integrable oscillation envelopes
\(\sup_{n\geq N}|Z_n-Z|\).
It supplies Exercise~4.7.2 after translating the source's
negative indices to \(n\to+\infty\).
The envelope cannot be replaced by uniform integrability, and an
\(L^1\) form holds under \(L^1\) convergence alone; see
\cref{thm:p15-w02b-p15-022-backward-changing-boundary}.
\end{hbtheorem}

\begin{proofboundary}{def:p15-exchangeability-source-scope}{The source applications and separate historical presentations}
The complete supplied Durrett chain is Examples~4.7.4--4.7.6 and 4.7.8,
Lemma~4.7.7, Theorems~4.7.9--4.7.10, the finite-urn discussion
and all five exercises. The earlier reverse convergence theorems
are \cref{thm:p15-backward-martingale-convergence,thm:p15-levy-backward}.
The regular-law and uniqueness construction states its standard
Borel domain; a regular probability-valued representation on every
arbitrary measurable space is not asserted.
Durrett's own boundary for arbitrary spaces (presentability, the
regular conditional law, and the Dubins--Freedman and Freedman
counterexample references) is in \cref{sec:p15-w02b-p15-022-presentable},
and his foundations-of-statistics remark follows
\cref{thm:p15-w02b-p15-022-predictive}.

Feller's separate completely-monotone moment presentation is supplied in
\cref{thm:p15-feller-binary-mixture}, with its full Bernstein and
Hausdorff prerequisites in \cref{sec:p15-feller-bernstein-hausdorff}
\cite[vol.~II, VII.1--VII.4]{Feller1971ProbabilityII}.
The full Diaconis--Freedman quantitative presentation and the source's
general-space counterexample references retain their own exact-source
obligations. The P\'olya beta mixture already proved in
\cref{thm:p15-polya-beta-limit} is one explicit instance of the
general binary representation here.
\end{proofboundary}
